\begin{textsample}{NEG Dim 8 – human – Score -1.00 – sp/sp\_ef\_matematica\_3\_p2.txt}  \label{ex:f8_neg_018}
As unidades temáticas reúnem um conjunto de ideias fundamentais, tais como : - Equivalência, presente nos estudos dos números racionais, equações, áreas ou volumes e em outros objetos de conhecimento ; - Ordem, está presente nos conjuntos numéricos, na construção de algoritmos e em outros procedimentos, como sequências e organização ; - Proporcionalidade, que contempla o raciocínio analógico, comparações quando se trata de frações, razões e proporções, semelhança de figuras, grandezas diretamente proporcionais, entre outros ; - Aproximação, que está articulada com a realização de cálculos aproximados, como estimativas e outros utilizados no dia a dia ; - Variação, conceito \textbf{associado} ao estudo das formas de crescimento e decrescimento, taxas de variação num dado contexto, como por \textbf{exemplo}, financeiro ; - Interdependência, \textbf{associada} à ideia de funções com ou sem uso de fórmulas, por \textbf{exemplo}, ligada à ideia de “ se p, então, q ”, sendo uma sentença matemática mais recorrente ; - Representação, \textbf{associada} à percepção e representação do espaço, de formas geométricas existentes ou imaginadas ; também \textbf{associada} aos números, às operações e à interdependência.

% matched lemmas: associar, exemplo
\end{textsample}
